\widowpenalty10000
\clubpenalty10000

\usepackage{tcolorbox}
\tcbuselibrary{skins}
\newtcolorbox{objectives}{
  standard jigsaw,
  opacityback=0,
  colframe=black,
  boxrule=1pt}

\newtcolorbox{Jeu}{
  enhanced,
  colback=white,
  colframe=white,
  boxrule=0pt,
  borderline west={2pt}{0pt}{orange},
  left=1em, right=0em, top=0.35em, bottom=0.3em}

\renewcommand{\arraystretch}{1.5}

\usepackage{multicol}

\usepackage{booktabs}
\usepackage{amsthm}
\usepackage{cancel}
\usepackage{gensymb}
\usepackage[version=4]{mhchem}
\usepackage{siunitx}
\sisetup{exponent-product = \cdot}
\usepackage{tabto}

\usepackage[answerdelayed,lastexercise]{exercise}
\renewcommand{\ExerciseName}{Exercice}
\renewcommand{\AnswerName}{Solution de l'exercice}

% Empêche que le titre « Solution de l'exercice N » soit séparé de sa résolution
% par un saut de page : on réserve quelques lignes avant le titre (sinon la
% coupure se fait avant lui) et on interdit la coupure juste après.
\usepackage{needspace}
\renewcommand{\AnswerHeader}{%
  \needspace{5\baselineskip}%
  \medskip\centerline{\textbf{\AnswerName\ \ExerciseHeaderNB}\smallskip}%
  \nobreak}

% De même pour le titre « Exercice N » : on le garde avec son énoncé.
% Ici le \medskip est APRÈS le titre, donc on place le \nobreak juste avant lui
% pour neutraliser la coupure possible entre le titre et l'énoncé.
\renewcommand{\ExerciseHeader}{%
  \needspace{5\baselineskip}%
  \centerline{\textbf{\large\ExerciseHeaderDifficulty\ExerciseName\ %
    \ExerciseHeaderNB\ExerciseHeaderTitle\ExerciseHeaderOrigin}}%
  \nobreak\medskip}

\usepackage{fancyhdr}
\setlength{\headheight}{15pt}
\pagestyle{fancy}
\lhead{}
\chead{}
\rhead{\leftmark}
\lfoot{MRzl - \the\year}
\cfoot{}
\rfoot{\thepage}
\renewcommand{\headrulewidth}{0.2pt}
\renewcommand{\footrulewidth}{0.2pt}
\renewcommand{\chaptermark}[1]{\markboth{\thechapter.\space#1}{}}

\makeatletter
\def\thm@space@setup{%
  \thm@preskip=8pt plus 2pt minus 4pt
  \thm@postskip=\thm@preskip
}
\makeatother

\newenvironment{credit}
  {\vspace{-1em}\begin{center}\begin{footnotesize}\begin{textit}}
  {\end{textit}\end{footnotesize}\end{center}}
